% Options for packages loaded elsewhere
\PassOptionsToPackage{unicode}{hyperref}
\PassOptionsToPackage{hyphens}{url}
\PassOptionsToPackage{dvipsnames,svgnames,x11names}{xcolor}
%
\documentclass[
  letterpaper,
  DIV=11,
  numbers=noendperiod]{scrreprt}

\usepackage{amsmath,amssymb}
\usepackage{iftex}
\ifPDFTeX
  \usepackage[T1]{fontenc}
  \usepackage[utf8]{inputenc}
  \usepackage{textcomp} % provide euro and other symbols
\else % if luatex or xetex
  \usepackage{unicode-math}
  \defaultfontfeatures{Scale=MatchLowercase}
  \defaultfontfeatures[\rmfamily]{Ligatures=TeX,Scale=1}
\fi
\usepackage{lmodern}
\ifPDFTeX\else  
    % xetex/luatex font selection
\fi
% Use upquote if available, for straight quotes in verbatim environments
\IfFileExists{upquote.sty}{\usepackage{upquote}}{}
\IfFileExists{microtype.sty}{% use microtype if available
  \usepackage[]{microtype}
  \UseMicrotypeSet[protrusion]{basicmath} % disable protrusion for tt fonts
}{}
\makeatletter
\@ifundefined{KOMAClassName}{% if non-KOMA class
  \IfFileExists{parskip.sty}{%
    \usepackage{parskip}
  }{% else
    \setlength{\parindent}{0pt}
    \setlength{\parskip}{6pt plus 2pt minus 1pt}}
}{% if KOMA class
  \KOMAoptions{parskip=half}}
\makeatother
\usepackage{xcolor}
\setlength{\emergencystretch}{3em} % prevent overfull lines
\setcounter{secnumdepth}{5}
% Make \paragraph and \subparagraph free-standing
\makeatletter
\ifx\paragraph\undefined\else
  \let\oldparagraph\paragraph
  \renewcommand{\paragraph}{
    \@ifstar
      \xxxParagraphStar
      \xxxParagraphNoStar
  }
  \newcommand{\xxxParagraphStar}[1]{\oldparagraph*{#1}\mbox{}}
  \newcommand{\xxxParagraphNoStar}[1]{\oldparagraph{#1}\mbox{}}
\fi
\ifx\subparagraph\undefined\else
  \let\oldsubparagraph\subparagraph
  \renewcommand{\subparagraph}{
    \@ifstar
      \xxxSubParagraphStar
      \xxxSubParagraphNoStar
  }
  \newcommand{\xxxSubParagraphStar}[1]{\oldsubparagraph*{#1}\mbox{}}
  \newcommand{\xxxSubParagraphNoStar}[1]{\oldsubparagraph{#1}\mbox{}}
\fi
\makeatother

\usepackage{color}
\usepackage{fancyvrb}
\newcommand{\VerbBar}{|}
\newcommand{\VERB}{\Verb[commandchars=\\\{\}]}
\DefineVerbatimEnvironment{Highlighting}{Verbatim}{commandchars=\\\{\}}
% Add ',fontsize=\small' for more characters per line
\usepackage{framed}
\definecolor{shadecolor}{RGB}{241,243,245}
\newenvironment{Shaded}{\begin{snugshade}}{\end{snugshade}}
\newcommand{\AlertTok}[1]{\textcolor[rgb]{0.68,0.00,0.00}{#1}}
\newcommand{\AnnotationTok}[1]{\textcolor[rgb]{0.37,0.37,0.37}{#1}}
\newcommand{\AttributeTok}[1]{\textcolor[rgb]{0.40,0.45,0.13}{#1}}
\newcommand{\BaseNTok}[1]{\textcolor[rgb]{0.68,0.00,0.00}{#1}}
\newcommand{\BuiltInTok}[1]{\textcolor[rgb]{0.00,0.23,0.31}{#1}}
\newcommand{\CharTok}[1]{\textcolor[rgb]{0.13,0.47,0.30}{#1}}
\newcommand{\CommentTok}[1]{\textcolor[rgb]{0.37,0.37,0.37}{#1}}
\newcommand{\CommentVarTok}[1]{\textcolor[rgb]{0.37,0.37,0.37}{\textit{#1}}}
\newcommand{\ConstantTok}[1]{\textcolor[rgb]{0.56,0.35,0.01}{#1}}
\newcommand{\ControlFlowTok}[1]{\textcolor[rgb]{0.00,0.23,0.31}{\textbf{#1}}}
\newcommand{\DataTypeTok}[1]{\textcolor[rgb]{0.68,0.00,0.00}{#1}}
\newcommand{\DecValTok}[1]{\textcolor[rgb]{0.68,0.00,0.00}{#1}}
\newcommand{\DocumentationTok}[1]{\textcolor[rgb]{0.37,0.37,0.37}{\textit{#1}}}
\newcommand{\ErrorTok}[1]{\textcolor[rgb]{0.68,0.00,0.00}{#1}}
\newcommand{\ExtensionTok}[1]{\textcolor[rgb]{0.00,0.23,0.31}{#1}}
\newcommand{\FloatTok}[1]{\textcolor[rgb]{0.68,0.00,0.00}{#1}}
\newcommand{\FunctionTok}[1]{\textcolor[rgb]{0.28,0.35,0.67}{#1}}
\newcommand{\ImportTok}[1]{\textcolor[rgb]{0.00,0.46,0.62}{#1}}
\newcommand{\InformationTok}[1]{\textcolor[rgb]{0.37,0.37,0.37}{#1}}
\newcommand{\KeywordTok}[1]{\textcolor[rgb]{0.00,0.23,0.31}{\textbf{#1}}}
\newcommand{\NormalTok}[1]{\textcolor[rgb]{0.00,0.23,0.31}{#1}}
\newcommand{\OperatorTok}[1]{\textcolor[rgb]{0.37,0.37,0.37}{#1}}
\newcommand{\OtherTok}[1]{\textcolor[rgb]{0.00,0.23,0.31}{#1}}
\newcommand{\PreprocessorTok}[1]{\textcolor[rgb]{0.68,0.00,0.00}{#1}}
\newcommand{\RegionMarkerTok}[1]{\textcolor[rgb]{0.00,0.23,0.31}{#1}}
\newcommand{\SpecialCharTok}[1]{\textcolor[rgb]{0.37,0.37,0.37}{#1}}
\newcommand{\SpecialStringTok}[1]{\textcolor[rgb]{0.13,0.47,0.30}{#1}}
\newcommand{\StringTok}[1]{\textcolor[rgb]{0.13,0.47,0.30}{#1}}
\newcommand{\VariableTok}[1]{\textcolor[rgb]{0.07,0.07,0.07}{#1}}
\newcommand{\VerbatimStringTok}[1]{\textcolor[rgb]{0.13,0.47,0.30}{#1}}
\newcommand{\WarningTok}[1]{\textcolor[rgb]{0.37,0.37,0.37}{\textit{#1}}}

\providecommand{\tightlist}{%
  \setlength{\itemsep}{0pt}\setlength{\parskip}{0pt}}\usepackage{longtable,booktabs,array}
\usepackage{calc} % for calculating minipage widths
% Correct order of tables after \paragraph or \subparagraph
\usepackage{etoolbox}
\makeatletter
\patchcmd\longtable{\par}{\if@noskipsec\mbox{}\fi\par}{}{}
\makeatother
% Allow footnotes in longtable head/foot
\IfFileExists{footnotehyper.sty}{\usepackage{footnotehyper}}{\usepackage{footnote}}
\makesavenoteenv{longtable}
\usepackage{graphicx}
\makeatletter
\newsavebox\pandoc@box
\newcommand*\pandocbounded[1]{% scales image to fit in text height/width
  \sbox\pandoc@box{#1}%
  \Gscale@div\@tempa{\textheight}{\dimexpr\ht\pandoc@box+\dp\pandoc@box\relax}%
  \Gscale@div\@tempb{\linewidth}{\wd\pandoc@box}%
  \ifdim\@tempb\p@<\@tempa\p@\let\@tempa\@tempb\fi% select the smaller of both
  \ifdim\@tempa\p@<\p@\scalebox{\@tempa}{\usebox\pandoc@box}%
  \else\usebox{\pandoc@box}%
  \fi%
}
% Set default figure placement to htbp
\def\fps@figure{htbp}
\makeatother
% definitions for citeproc citations
\NewDocumentCommand\citeproctext{}{}
\NewDocumentCommand\citeproc{mm}{%
  \begingroup\def\citeproctext{#2}\cite{#1}\endgroup}
\makeatletter
 % allow citations to break across lines
 \let\@cite@ofmt\@firstofone
 % avoid brackets around text for \cite:
 \def\@biblabel#1{}
 \def\@cite#1#2{{#1\if@tempswa , #2\fi}}
\makeatother
\newlength{\cslhangindent}
\setlength{\cslhangindent}{1.5em}
\newlength{\csllabelwidth}
\setlength{\csllabelwidth}{3em}
\newenvironment{CSLReferences}[2] % #1 hanging-indent, #2 entry-spacing
 {\begin{list}{}{%
  \setlength{\itemindent}{0pt}
  \setlength{\leftmargin}{0pt}
  \setlength{\parsep}{0pt}
  % turn on hanging indent if param 1 is 1
  \ifodd #1
   \setlength{\leftmargin}{\cslhangindent}
   \setlength{\itemindent}{-1\cslhangindent}
  \fi
  % set entry spacing
  \setlength{\itemsep}{#2\baselineskip}}}
 {\end{list}}
\usepackage{calc}
\newcommand{\CSLBlock}[1]{\hfill\break\parbox[t]{\linewidth}{\strut\ignorespaces#1\strut}}
\newcommand{\CSLLeftMargin}[1]{\parbox[t]{\csllabelwidth}{\strut#1\strut}}
\newcommand{\CSLRightInline}[1]{\parbox[t]{\linewidth - \csllabelwidth}{\strut#1\strut}}
\newcommand{\CSLIndent}[1]{\hspace{\cslhangindent}#1}

\KOMAoption{captions}{tableheading}
\makeatletter
\@ifpackageloaded{tcolorbox}{}{\usepackage[skins,breakable]{tcolorbox}}
\@ifpackageloaded{fontawesome5}{}{\usepackage{fontawesome5}}
\definecolor{quarto-callout-color}{HTML}{909090}
\definecolor{quarto-callout-note-color}{HTML}{0758E5}
\definecolor{quarto-callout-important-color}{HTML}{CC1914}
\definecolor{quarto-callout-warning-color}{HTML}{EB9113}
\definecolor{quarto-callout-tip-color}{HTML}{00A047}
\definecolor{quarto-callout-caution-color}{HTML}{FC5300}
\definecolor{quarto-callout-color-frame}{HTML}{acacac}
\definecolor{quarto-callout-note-color-frame}{HTML}{4582ec}
\definecolor{quarto-callout-important-color-frame}{HTML}{d9534f}
\definecolor{quarto-callout-warning-color-frame}{HTML}{f0ad4e}
\definecolor{quarto-callout-tip-color-frame}{HTML}{02b875}
\definecolor{quarto-callout-caution-color-frame}{HTML}{fd7e14}
\makeatother
\makeatletter
\@ifpackageloaded{bookmark}{}{\usepackage{bookmark}}
\makeatother
\makeatletter
\@ifpackageloaded{caption}{}{\usepackage{caption}}
\AtBeginDocument{%
\ifdefined\contentsname
  \renewcommand*\contentsname{Table of contents}
\else
  \newcommand\contentsname{Table of contents}
\fi
\ifdefined\listfigurename
  \renewcommand*\listfigurename{List of Figures}
\else
  \newcommand\listfigurename{List of Figures}
\fi
\ifdefined\listtablename
  \renewcommand*\listtablename{List of Tables}
\else
  \newcommand\listtablename{List of Tables}
\fi
\ifdefined\figurename
  \renewcommand*\figurename{Figure}
\else
  \newcommand\figurename{Figure}
\fi
\ifdefined\tablename
  \renewcommand*\tablename{Table}
\else
  \newcommand\tablename{Table}
\fi
}
\@ifpackageloaded{float}{}{\usepackage{float}}
\floatstyle{ruled}
\@ifundefined{c@chapter}{\newfloat{codelisting}{h}{lop}}{\newfloat{codelisting}{h}{lop}[chapter]}
\floatname{codelisting}{Listing}
\newcommand*\listoflistings{\listof{codelisting}{List of Listings}}
\makeatother
\makeatletter
\makeatother
\makeatletter
\@ifpackageloaded{caption}{}{\usepackage{caption}}
\@ifpackageloaded{subcaption}{}{\usepackage{subcaption}}
\makeatother

\usepackage{bookmark}

\IfFileExists{xurl.sty}{\usepackage{xurl}}{} % add URL line breaks if available
\urlstyle{same} % disable monospaced font for URLs
\hypersetup{
  pdftitle={Evaluating Diagnostic Tests with ROC Curves in R},
  pdfauthor={Brynn Baumert},
  colorlinks=true,
  linkcolor={blue},
  filecolor={Maroon},
  citecolor={Blue},
  urlcolor={Blue},
  pdfcreator={LaTeX via pandoc}}


\title{Evaluating Diagnostic Tests with ROC Curves in R}
\author{Brynn Baumert}
\date{March 6, 2026}

\begin{document}
\maketitle

\renewcommand*\contentsname{Table of contents}
{
\hypersetup{linkcolor=}
\setcounter{tocdepth}{2}
\tableofcontents
}

\bookmarksetup{startatroot}

\chapter*{Assignment Flow}\label{assignment-flow}
\addcontentsline{toc}{chapter}{Assignment Flow}

\markboth{Assignment Flow}{Assignment Flow}

This project is set up to scaffold this user guide and set you up for
success.

Start at \texttt{topic.qmd} and identify a topic for your user guide.

Once you've proposed a topic and received my approval/feedback, you can
convert the contents of \texttt{topic.qmd} into an introductory
paragraph in your report. Please leave the \texttt{topic.qmd} file as an
appendix, to document the different stages of this project.

Once we've agreed on a topic for your user guide, you will proceed to
\texttt{needs.qmd} and \texttt{task-analysis.qmd}. Both will be
submitted at the same check-in, and both will be included in your final
report as appendices. You should be able to re-purpose most of the
information in \texttt{task-analysis.qmd} to orient the user to
different components of the task you've chosen.

Your next step is to actually write the content that satisfies the work
you outlined in \texttt{task-analysis.qmd}.

Once you're finished with that, ideally, you'll have plenty of time to
edit and streamline your report. Feel free to trade reports with a
friend and try to complete your friend's task - this will help you both
identify areas where the guide isn't as clear. Try to complete the task
with a different data set - that often helps find trouble spots.

When you submit your final report, you may remove \texttt{index.qmd}
from the \texttt{\_quarto.yml} file, which will remove this chapter
(which is quite unnecessary for your user) from the report. You should
also take the time to make sure your name is listed as both author and
copyright holder in the \texttt{\_quarto.yml} file. Feel free to make a
custom cover for your manual if you would like to do so. You can also
tweak the CSS/theme for the book, so long as you're conscious of
accessibility concerns.

\begin{tcolorbox}[enhanced jigsaw, opacityback=0, opacitybacktitle=0.6, toptitle=1mm, leftrule=.75mm, colback=white, titlerule=0mm, title=\textcolor{quarto-callout-caution-color}{\faFire}\hspace{0.5em}{Building the report}, coltitle=black, colbacktitle=quarto-callout-caution-color!10!white, left=2mm, bottomrule=.15mm, bottomtitle=1mm, arc=.35mm, toprule=.15mm, rightrule=.15mm, breakable, colframe=quarto-callout-caution-color-frame]

If you are using RStudio to complete this report, you can hit
(Ctrl/Cmd)-Shift-B to build the whole report. You can also type
\texttt{quarto\ render\ .} on the command line in the project folder to
accomplish the same thing.

If you have questions about how to customize or debug your book, it may
be helpful to start at the quarto book documentation
(\textbf{positpbcCreatingBook2024?}).

\end{tcolorbox}

\bookmarksetup{startatroot}

\chapter{Evaluating Diagnostic Tests with ROC Curves in
R}\label{evaluating-diagnostic-tests-with-roc-curves-in-r}

\section{Introduction}\label{introduction}

Diagnostic tests in medicine and biology rarely produce a simple yes/no
result. More often, they produce a continuous measurement, for example a
biomarker concentration, an imaging score, a blood pressure reading. The
researcher must decide how well that measurement distinguishes between
individuals with and without a condition. Receiver Operating
Characteristic (ROC) curves are one of the most widely used tools for
evaluating this kind of diagnostic performance.

This guide walks through the full ROC analysis workflow in R: deciding
whether ROC analysis is appropriate for your data, preparing your
dataset, generating and visualizing the ROC curve, and computing and
interpreting the Area Under the Curve (AUC). All steps are demonstrated
with real datasets and annotated code using the \texttt{pROC} package
(Robin et al. 2011).

This guide does \textbf{not} cover advanced topics such as
covariate-adjusted ROC curves, time-dependent ROC analysis, or optimal
threshold selection (e.g., the Youden Index). Those topics are excluded
to keep the focus on foundational ROC analysis.

\begin{center}\rule{0.5\linewidth}{0.5pt}\end{center}

\section{Prerequisites}\label{prerequisites}

\subsection{Knowledge}\label{knowledge}

You should be familiar with:

\begin{itemize}
\tightlist
\item
  \textbf{Sensitivity}: the probability that the test is positive given
  the person has the disease (true positive rate)
\item
  \textbf{Specificity}: the probability that the test is negative given
  the person does not have the disease (true negative rate)
\item
  \textbf{Binary outcomes}: your reference standard should classify each
  individual as diseased (1) or non-diseased (0)
\end{itemize}

No prior experience with ROC curves is needed

\subsection{Software}\label{software}

You will need R and RStudio installed, along with the following
packages:

\begin{Shaded}
\begin{Highlighting}[]
\CommentTok{\#loading the following packages and install if necessary}
\FunctionTok{library}\NormalTok{(pROC)}
\end{Highlighting}
\end{Shaded}

\begin{verbatim}
Warning: package 'pROC' was built under R version 4.4.3
\end{verbatim}

\begin{verbatim}
Type 'citation("pROC")' for a citation.
\end{verbatim}

\begin{verbatim}

Attaching package: 'pROC'
\end{verbatim}

\begin{verbatim}
The following objects are masked from 'package:stats':

    cov, smooth, var
\end{verbatim}

\begin{Shaded}
\begin{Highlighting}[]
\FunctionTok{library}\NormalTok{(ggplot2)}
\end{Highlighting}
\end{Shaded}

\begin{verbatim}
Warning: package 'ggplot2' was built under R version 4.4.3
\end{verbatim}

\begin{Shaded}
\begin{Highlighting}[]
\FunctionTok{library}\NormalTok{(mlbench)}
\end{Highlighting}
\end{Shaded}

\begin{verbatim}
Warning: package 'mlbench' was built under R version 4.4.3
\end{verbatim}

\section{The Data}\label{the-data}

This guide uses three datasets to demonstrate ROC analysis across
different applied contexts.

\subsection{\texorpdfstring{1. \texttt{aSAH} --- Aneurysmal Subarachnoid
Hemorrhage}{1. aSAH --- Aneurysmal Subarachnoid Hemorrhage}}\label{asah-aneurysmal-subarachnoid-hemorrhage}

Built into the \texttt{pROC} package. The dataset contains clinical and
biomarker data from patients with subarachnoid hemorrhage. We use the
biomarker \texttt{s100b} to predict a poor neurological outcome
(\texttt{outcome}).

\begin{Shaded}
\begin{Highlighting}[]
\CommentTok{\#load and preview dataset{-}asah}

\FunctionTok{data}\NormalTok{(aSAH)}
\FunctionTok{head}\NormalTok{(aSAH)}
\end{Highlighting}
\end{Shaded}

\begin{verbatim}
   gos6 outcome gender age wfns s100b  ndka
29    5    Good Female  42    1  0.13  3.01
30    5    Good Female  37    1  0.14  8.54
31    5    Good Female  42    1  0.10  8.09
32    5    Good Female  27    1  0.04 10.42
33    1    Poor Female  42    3  0.13 17.40
34    1    Poor   Male  48    2  0.10 12.75
\end{verbatim}

\begin{Shaded}
\begin{Highlighting}[]
\FunctionTok{str}\NormalTok{(aSAH)}
\end{Highlighting}
\end{Shaded}

\begin{verbatim}
'data.frame':   113 obs. of  7 variables:
 $ gos6   : Ord.factor w/ 5 levels "1"<"2"<"3"<"4"<..: 5 5 5 5 1 1 4 1 5 4 ...
 $ outcome: Factor w/ 2 levels "Good","Poor": 1 1 1 1 2 2 1 2 1 1 ...
 $ gender : Factor w/ 2 levels "Male","Female": 2 2 2 2 2 1 1 1 2 2 ...
 $ age    : int  42 37 42 27 42 48 57 41 49 75 ...
 $ wfns   : Ord.factor w/ 5 levels "1"<"2"<"3"<"4"<..: 1 1 1 1 3 2 5 4 1 2 ...
 $ s100b  : num  0.13 0.14 0.1 0.04 0.13 0.1 0.47 0.16 0.18 0.1 ...
 $ ndka   : num  3.01 8.54 8.09 10.42 17.4 ...
\end{verbatim}

\subsection{\texorpdfstring{2. \texttt{PimaIndiansDiabetes} --- Diabetes
Screening}{2. PimaIndiansDiabetes --- Diabetes Screening}}\label{pimaindiansdiabetes-diabetes-screening}

From the \texttt{mlbench} package, the dataset contains health
measurements for Pima Indian women. We use \texttt{glucose} to predict
diabetes status (\texttt{diabetes}).

\begin{Shaded}
\begin{Highlighting}[]
\CommentTok{\#load and preview data set{-}pima}

\FunctionTok{data}\NormalTok{(PimaIndiansDiabetes)}
\FunctionTok{head}\NormalTok{(PimaIndiansDiabetes)}
\end{Highlighting}
\end{Shaded}

\begin{verbatim}
  pregnant glucose pressure triceps insulin mass pedigree age diabetes
1        6     148       72      35       0 33.6    0.627  50      pos
2        1      85       66      29       0 26.6    0.351  31      neg
3        8     183       64       0       0 23.3    0.672  32      pos
4        1      89       66      23      94 28.1    0.167  21      neg
5        0     137       40      35     168 43.1    2.288  33      pos
6        5     116       74       0       0 25.6    0.201  30      neg
\end{verbatim}

\begin{Shaded}
\begin{Highlighting}[]
\FunctionTok{str}\NormalTok{(PimaIndiansDiabetes)}
\end{Highlighting}
\end{Shaded}

\begin{verbatim}
'data.frame':   768 obs. of  9 variables:
 $ pregnant: num  6 1 8 1 0 5 3 10 2 8 ...
 $ glucose : num  148 85 183 89 137 116 78 115 197 125 ...
 $ pressure: num  72 66 64 66 40 74 50 0 70 96 ...
 $ triceps : num  35 29 0 23 35 0 32 0 45 0 ...
 $ insulin : num  0 0 0 94 168 0 88 0 543 0 ...
 $ mass    : num  33.6 26.6 23.3 28.1 43.1 25.6 31 35.3 30.5 0 ...
 $ pedigree: num  0.627 0.351 0.672 0.167 2.288 ...
 $ age     : num  50 31 32 21 33 30 26 29 53 54 ...
 $ diabetes: Factor w/ 2 levels "neg","pos": 2 1 2 1 2 1 2 1 2 2 ...
\end{verbatim}

\subsection{3. Cleveland Heart Disease
Dataset}\label{cleveland-heart-disease-dataset}

Loaded directly from the UCI Machine Learning Repository. We use
\texttt{thalach} (maximum heart rate achieved) to predict heart disease
presence (\texttt{target}).

\begin{Shaded}
\begin{Highlighting}[]
\CommentTok{\#load{-}heart}
\NormalTok{heart }\OtherTok{\textless{}{-}} \FunctionTok{read.csv}\NormalTok{(}
  \StringTok{"https://archive.ics.uci.edu/ml/machine{-}learning{-}databases/heart{-}disease/processed.cleveland.data"}\NormalTok{,}
  \AttributeTok{header =} \ConstantTok{FALSE}\NormalTok{,}
  \AttributeTok{na.strings =} \StringTok{"?"}
\NormalTok{)}
\FunctionTok{colnames}\NormalTok{(heart) }\OtherTok{\textless{}{-}} \FunctionTok{c}\NormalTok{(}\StringTok{"age"}\NormalTok{, }\StringTok{"sex"}\NormalTok{, }\StringTok{"cp"}\NormalTok{, }\StringTok{"trestbps"}\NormalTok{, }\StringTok{"chol"}\NormalTok{,}
                     \StringTok{"fbs"}\NormalTok{, }\StringTok{"restecg"}\NormalTok{, }\StringTok{"thalach"}\NormalTok{, }\StringTok{"exang"}\NormalTok{,}
                     \StringTok{"oldpeak"}\NormalTok{, }\StringTok{"slope"}\NormalTok{, }\StringTok{"ca"}\NormalTok{, }\StringTok{"thal"}\NormalTok{, }\StringTok{"target"}\NormalTok{)}
\NormalTok{heart}\SpecialCharTok{$}\NormalTok{target }\OtherTok{\textless{}{-}} \FunctionTok{ifelse}\NormalTok{(heart}\SpecialCharTok{$}\NormalTok{target }\SpecialCharTok{\textgreater{}} \DecValTok{0}\NormalTok{, }\DecValTok{1}\NormalTok{, }\DecValTok{0}\NormalTok{)}
\FunctionTok{head}\NormalTok{(heart)}
\end{Highlighting}
\end{Shaded}

\begin{verbatim}
  age sex cp trestbps chol fbs restecg thalach exang oldpeak slope ca thal
1  63   1  1      145  233   1       2     150     0     2.3     3  0    6
2  67   1  4      160  286   0       2     108     1     1.5     2  3    3
3  67   1  4      120  229   0       2     129     1     2.6     2  2    7
4  37   1  3      130  250   0       0     187     0     3.5     3  0    3
5  41   0  2      130  204   0       2     172     0     1.4     1  0    3
6  56   1  2      120  236   0       0     178     0     0.8     1  0    3
  target
1      0
2      1
3      1
4      0
5      0
6      0
\end{verbatim}

\begin{Shaded}
\begin{Highlighting}[]
\FunctionTok{str}\NormalTok{(heart)}
\end{Highlighting}
\end{Shaded}

\begin{verbatim}
'data.frame':   303 obs. of  14 variables:
 $ age     : num  63 67 67 37 41 56 62 57 63 53 ...
 $ sex     : num  1 1 1 1 0 1 0 0 1 1 ...
 $ cp      : num  1 4 4 3 2 2 4 4 4 4 ...
 $ trestbps: num  145 160 120 130 130 120 140 120 130 140 ...
 $ chol    : num  233 286 229 250 204 236 268 354 254 203 ...
 $ fbs     : num  1 0 0 0 0 0 0 0 0 1 ...
 $ restecg : num  2 2 2 0 2 0 2 0 2 2 ...
 $ thalach : num  150 108 129 187 172 178 160 163 147 155 ...
 $ exang   : num  0 1 1 0 0 0 0 1 0 1 ...
 $ oldpeak : num  2.3 1.5 2.6 3.5 1.4 0.8 3.6 0.6 1.4 3.1 ...
 $ slope   : num  3 2 2 3 1 1 3 1 2 3 ...
 $ ca      : num  0 3 2 0 0 0 2 0 1 0 ...
 $ thal    : num  6 3 7 3 3 3 3 3 7 7 ...
 $ target  : num  0 1 1 0 0 0 1 0 1 1 ...
\end{verbatim}

\section{Step 1: Is ROC Analysis Appropriate for Your
Data?}\label{step-1-is-roc-analysis-appropriate-for-your-data}

Before writing any code, confirm that ROC analysis is the right tool.

\textbf{ROC analysis is appropriate when:}

\begin{itemize}
\tightlist
\item
  Your diagnostic test produces a \textbf{continuous or ordinal} result
  (e.g., a blood glucose level, a probability score)
\item
  You have a \textbf{binary reference standard} --- each observation is
  either diseased (1) or non-diseased (0)
\item
  Your goal is to evaluate \textbf{overall discriminative ability}
  across all possible thresholds, not performance at one specific cutoff
\end{itemize}

\textbf{ROC analysis is not appropriate when:}

\begin{itemize}
\tightlist
\item
  Your test result is already \textbf{binary} (positive/negative)
\item
  Disease \textbf{prevalence is very low}
\item
  The \textbf{costs of false positives and false negatives are very
  different} --- ROC curves treat both error types equally
\end{itemize}

\subsection{Quick Checklist}\label{quick-checklist}

\begin{longtable}[]{@{}ll@{}}
\toprule\noalign{}
Question & Required Answer \\
\midrule\noalign{}
\endhead
\bottomrule\noalign{}
\endlastfoot
Is the test result continuous or ordinal? & Yes \\
Is the outcome binary (diseased / not diseased)? & Yes \\
Am I interested in performance across all thresholds? & Yes \\
Is disease prevalence reasonably balanced? & Ideally yes \\
\end{longtable}

\section{Step 2: Prepare Your Data}\label{step-2-prepare-your-data}

Your dataset must have at minimum:

\begin{itemize}
\tightlist
\item
  A \textbf{numeric column} for the test result
\item
  A \textbf{binary column} for the outcome (0/1 or a two-level factor)
\end{itemize}

Check your data structure and look for missing values before proceeding.

\begin{Shaded}
\begin{Highlighting}[]
\CommentTok{\# Check structure}
\FunctionTok{str}\NormalTok{(aSAH)}
\end{Highlighting}
\end{Shaded}

\begin{verbatim}
'data.frame':   113 obs. of  7 variables:
 $ gos6   : Ord.factor w/ 5 levels "1"<"2"<"3"<"4"<..: 5 5 5 5 1 1 4 1 5 4 ...
 $ outcome: Factor w/ 2 levels "Good","Poor": 1 1 1 1 2 2 1 2 1 1 ...
 $ gender : Factor w/ 2 levels "Male","Female": 2 2 2 2 2 1 1 1 2 2 ...
 $ age    : int  42 37 42 27 42 48 57 41 49 75 ...
 $ wfns   : Ord.factor w/ 5 levels "1"<"2"<"3"<"4"<..: 1 1 1 1 3 2 5 4 1 2 ...
 $ s100b  : num  0.13 0.14 0.1 0.04 0.13 0.1 0.47 0.16 0.18 0.1 ...
 $ ndka   : num  3.01 8.54 8.09 10.42 17.4 ...
\end{verbatim}

\begin{Shaded}
\begin{Highlighting}[]
\CommentTok{\# Check for missing values}
\FunctionTok{sum}\NormalTok{(}\FunctionTok{is.na}\NormalTok{(aSAH}\SpecialCharTok{$}\NormalTok{s100b))}
\end{Highlighting}
\end{Shaded}

\begin{verbatim}
[1] 0
\end{verbatim}

\begin{Shaded}
\begin{Highlighting}[]
\FunctionTok{sum}\NormalTok{(}\FunctionTok{is.na}\NormalTok{(aSAH}\SpecialCharTok{$}\NormalTok{outcome))}
\end{Highlighting}
\end{Shaded}

\begin{verbatim}
[1] 0
\end{verbatim}

\begin{Shaded}
\begin{Highlighting}[]
\CommentTok{\# Confirm outcome levels}
\FunctionTok{table}\NormalTok{(aSAH}\SpecialCharTok{$}\NormalTok{outcome)}
\end{Highlighting}
\end{Shaded}

\begin{verbatim}

Good Poor 
  72   41 
\end{verbatim}

The \texttt{pROC} package drops missing values automatically, but you
should know how many observations are affected. If missingness is
substantial or non-random, address it before running the analysis.

\section{Step 3: Generate the ROC
Curve}\label{step-3-generate-the-roc-curve}

Use \texttt{pROC::roc()} to fit the ROC curve. The two key arguments
are:

\begin{itemize}
\tightlist
\item
  \texttt{response}: the binary outcome variable
\item
  \texttt{predictor}: the continuous test result
\end{itemize}

\begin{Shaded}
\begin{Highlighting}[]
\NormalTok{roc\_asah }\OtherTok{\textless{}{-}} \FunctionTok{roc}\NormalTok{(}\AttributeTok{response =}\NormalTok{ aSAH}\SpecialCharTok{$}\NormalTok{outcome, }\AttributeTok{predictor =}\NormalTok{ aSAH}\SpecialCharTok{$}\NormalTok{s100b)}
\end{Highlighting}
\end{Shaded}

\begin{verbatim}
Setting levels: control = Good, case = Poor
\end{verbatim}

\begin{verbatim}
Setting direction: controls < cases
\end{verbatim}

\begin{Shaded}
\begin{Highlighting}[]
\NormalTok{roc\_asah}
\end{Highlighting}
\end{Shaded}

\begin{verbatim}

Call:
roc.default(response = aSAH$outcome, predictor = aSAH$s100b)

Data: aSAH$s100b in 72 controls (aSAH$outcome Good) < 41 cases (aSAH$outcome Poor).
Area under the curve: 0.7314
\end{verbatim}

\subsection{Understanding the Direction
Warning}\label{understanding-the-direction-warning}

\texttt{pROC} will print a message about the \texttt{direction} it
assumed (e.g., \texttt{direction:\ controls\ \textless{}\ cases}). This
tells you whether higher test values are assumed to indicate disease.
\textbf{Do not ignore this.} If the assumed direction is wrong for your
test, your curve will be mirrored and your AUC will be below 0.5. You
can set the direction explicitly:

\begin{Shaded}
\begin{Highlighting}[]
\CommentTok{\# Higher s100b {-}\textgreater{} worse outcome: controls \textless{} cases (default is correct here)}
\NormalTok{roc\_asah }\OtherTok{\textless{}{-}} \FunctionTok{roc}\NormalTok{(}\AttributeTok{response =}\NormalTok{ aSAH}\SpecialCharTok{$}\NormalTok{outcome,}
                \AttributeTok{predictor =}\NormalTok{ aSAH}\SpecialCharTok{$}\NormalTok{s100b,}
                \AttributeTok{direction =} \StringTok{"\textless{}"}\NormalTok{)}
\end{Highlighting}
\end{Shaded}

\begin{verbatim}
Setting levels: control = Good, case = Poor
\end{verbatim}

\section{Step 4: Compute and Interpret
AUC}\label{step-4-compute-and-interpret-auc}

The Area Under the Curve (AUC) summarizes the ROC curve in a single
number.

\begin{itemize}
\tightlist
\item
  \textbf{AUC = 0.5}: the test performs no better than random chance
\item
  \textbf{AUC = 1.0}: the test perfectly separates diseased from
  non-diseased individuals
\item
  \textbf{AUC = 0.7--0.8}: generally considered acceptable
  discrimination
\item
  \textbf{AUC \textgreater{} 0.8}: generally considered excellent
  discrimination
\end{itemize}

These are rough benchmarks --- what counts as ``good enough'' depends on
your clinical or scientific context.

\begin{Shaded}
\begin{Highlighting}[]
\FunctionTok{auc}\NormalTok{(roc\_asah)}
\end{Highlighting}
\end{Shaded}

\begin{verbatim}
Area under the curve: 0.7314
\end{verbatim}

\begin{Shaded}
\begin{Highlighting}[]
\FunctionTok{ci.auc}\NormalTok{(roc\_asah)}
\end{Highlighting}
\end{Shaded}

\begin{verbatim}
95% CI: 0.6301-0.8326 (DeLong)
\end{verbatim}

Repeat for the other datasets:

\begin{Shaded}
\begin{Highlighting}[]
\NormalTok{roc\_pima }\OtherTok{\textless{}{-}} \FunctionTok{roc}\NormalTok{(}\AttributeTok{response =}\NormalTok{ PimaIndiansDiabetes}\SpecialCharTok{$}\NormalTok{diabetes,}
                \AttributeTok{predictor =}\NormalTok{ PimaIndiansDiabetes}\SpecialCharTok{$}\NormalTok{glucose,}
                \AttributeTok{direction =} \StringTok{"\textless{}"}\NormalTok{)}
\end{Highlighting}
\end{Shaded}

\begin{verbatim}
Setting levels: control = neg, case = pos
\end{verbatim}

\begin{Shaded}
\begin{Highlighting}[]
\FunctionTok{auc}\NormalTok{(roc\_pima)}
\end{Highlighting}
\end{Shaded}

\begin{verbatim}
Area under the curve: 0.7881
\end{verbatim}

\begin{Shaded}
\begin{Highlighting}[]
\FunctionTok{ci.auc}\NormalTok{(roc\_pima)}
\end{Highlighting}
\end{Shaded}

\begin{verbatim}
95% CI: 0.7546-0.8217 (DeLong)
\end{verbatim}

\begin{Shaded}
\begin{Highlighting}[]
\NormalTok{heart\_clean }\OtherTok{\textless{}{-}} \FunctionTok{na.omit}\NormalTok{(heart[, }\FunctionTok{c}\NormalTok{(}\StringTok{"thalach"}\NormalTok{, }\StringTok{"target"}\NormalTok{)])}
\NormalTok{roc\_heart }\OtherTok{\textless{}{-}} \FunctionTok{roc}\NormalTok{(}\AttributeTok{response =}\NormalTok{ heart\_clean}\SpecialCharTok{$}\NormalTok{target,}
                 \AttributeTok{predictor =}\NormalTok{ heart\_clean}\SpecialCharTok{$}\NormalTok{thalach,}
                 \AttributeTok{direction =} \StringTok{"\textgreater{}"}\NormalTok{)  }\CommentTok{\# higher heart rate {-}\textgreater{} less likely disease}
\end{Highlighting}
\end{Shaded}

\begin{verbatim}
Setting levels: control = 0, case = 1
\end{verbatim}

\begin{Shaded}
\begin{Highlighting}[]
\FunctionTok{auc}\NormalTok{(roc\_heart)}
\end{Highlighting}
\end{Shaded}

\begin{verbatim}
Area under the curve: 0.7453
\end{verbatim}

\begin{Shaded}
\begin{Highlighting}[]
\FunctionTok{ci.auc}\NormalTok{(roc\_heart)}
\end{Highlighting}
\end{Shaded}

\begin{verbatim}
95% CI: 0.6901-0.8005 (DeLong)
\end{verbatim}

\section{Step 5: Visualize the ROC
Curve}\label{step-5-visualize-the-roc-curve}

A good visualization shows the curve, the diagonal reference line
(random chance), and the AUC value.

\subsection{\texorpdfstring{Basic Plot with
\texttt{pROC}}{Basic Plot with pROC}}\label{basic-plot-with-proc}

\begin{Shaded}
\begin{Highlighting}[]
\FunctionTok{plot}\NormalTok{(roc\_asah,}
     \AttributeTok{main =} \StringTok{"ROC Curve: s100b predicting poor outcome (aSAH)"}\NormalTok{,}
     \AttributeTok{col =} \StringTok{"\#2c7bb6"}\NormalTok{,}
     \AttributeTok{lwd =} \DecValTok{2}\NormalTok{)}
\FunctionTok{legend}\NormalTok{(}\StringTok{"bottomright"}\NormalTok{,}
       \AttributeTok{legend =} \FunctionTok{paste}\NormalTok{(}\StringTok{"AUC ="}\NormalTok{, }\FunctionTok{round}\NormalTok{(}\FunctionTok{auc}\NormalTok{(roc\_asah), }\DecValTok{3}\NormalTok{)),}
       \AttributeTok{col =} \StringTok{"\#2c7bb6"}\NormalTok{, }\AttributeTok{lwd =} \DecValTok{2}\NormalTok{)}
\end{Highlighting}
\end{Shaded}

\pandocbounded{\includegraphics[keepaspectratio]{guide_files/figure-pdf/plot-basic-1.pdf}}

\subsection{\texorpdfstring{Polished Plot with
\texttt{ggplot2}}{Polished Plot with ggplot2}}\label{polished-plot-with-ggplot2}

\begin{Shaded}
\begin{Highlighting}[]
\NormalTok{roc\_df }\OtherTok{\textless{}{-}} \FunctionTok{data.frame}\NormalTok{(}
  \AttributeTok{specificity =}\NormalTok{ roc\_asah}\SpecialCharTok{$}\NormalTok{specificities,}
  \AttributeTok{sensitivity =}\NormalTok{ roc\_asah}\SpecialCharTok{$}\NormalTok{sensitivities}
\NormalTok{)}

\FunctionTok{ggplot}\NormalTok{(roc\_df, }\FunctionTok{aes}\NormalTok{(}\AttributeTok{x =} \DecValTok{1} \SpecialCharTok{{-}}\NormalTok{ specificity, }\AttributeTok{y =}\NormalTok{ sensitivity)) }\SpecialCharTok{+}
  \FunctionTok{geom\_line}\NormalTok{(}\AttributeTok{color =} \StringTok{"\#2c7bb6"}\NormalTok{, }\AttributeTok{linewidth =} \FloatTok{1.2}\NormalTok{) }\SpecialCharTok{+}
  \FunctionTok{geom\_abline}\NormalTok{(}\AttributeTok{slope =} \DecValTok{1}\NormalTok{, }\AttributeTok{intercept =} \DecValTok{0}\NormalTok{, }\AttributeTok{linetype =} \StringTok{"dashed"}\NormalTok{, }\AttributeTok{color =} \StringTok{"gray50"}\NormalTok{) }\SpecialCharTok{+}
  \FunctionTok{annotate}\NormalTok{(}\StringTok{"text"}\NormalTok{, }\AttributeTok{x =} \FloatTok{0.75}\NormalTok{, }\AttributeTok{y =} \FloatTok{0.1}\NormalTok{,}
           \AttributeTok{label =} \FunctionTok{paste}\NormalTok{(}\StringTok{"AUC ="}\NormalTok{, }\FunctionTok{round}\NormalTok{(}\FunctionTok{auc}\NormalTok{(roc\_asah), }\DecValTok{3}\NormalTok{)),}
           \AttributeTok{size =} \DecValTok{5}\NormalTok{, }\AttributeTok{color =} \StringTok{"\#2c7bb6"}\NormalTok{) }\SpecialCharTok{+}
  \FunctionTok{labs}\NormalTok{(}\AttributeTok{title =} \StringTok{"ROC Curve: s100b predicting poor outcome (aSAH)"}\NormalTok{,}
       \AttributeTok{x =} \StringTok{"1 {-} Specificity (False Positive Rate)"}\NormalTok{,}
       \AttributeTok{y =} \StringTok{"Sensitivity (True Positive Rate)"}\NormalTok{) }\SpecialCharTok{+}
  \FunctionTok{theme\_minimal}\NormalTok{()}
\end{Highlighting}
\end{Shaded}

\pandocbounded{\includegraphics[keepaspectratio]{guide_files/figure-pdf/plot-ggplot-1.pdf}}

\subsection{Comparing Multiple ROC
Curves}\label{comparing-multiple-roc-curves}

\begin{Shaded}
\begin{Highlighting}[]
\FunctionTok{ggroc}\NormalTok{(}\FunctionTok{list}\NormalTok{(}\AttributeTok{aSAH =}\NormalTok{ roc\_asah, }\AttributeTok{Diabetes =}\NormalTok{ roc\_pima, }\AttributeTok{Heart =}\NormalTok{ roc\_heart),}
      \AttributeTok{linewidth =} \FloatTok{1.1}\NormalTok{) }\SpecialCharTok{+}
  \FunctionTok{geom\_abline}\NormalTok{(}\AttributeTok{slope =} \DecValTok{1}\NormalTok{, }\AttributeTok{intercept =} \DecValTok{1}\NormalTok{, }\AttributeTok{linetype =} \StringTok{"dashed"}\NormalTok{, }\AttributeTok{color =} \StringTok{"gray50"}\NormalTok{) }\SpecialCharTok{+}
  \FunctionTok{labs}\NormalTok{(}\AttributeTok{title =} \StringTok{"ROC Curve Comparison Across Datasets"}\NormalTok{,}
       \AttributeTok{x =} \StringTok{"1 {-} Specificity"}\NormalTok{, }\AttributeTok{y =} \StringTok{"Sensitivity"}\NormalTok{,}
       \AttributeTok{color =} \StringTok{"Dataset"}\NormalTok{) }\SpecialCharTok{+}
  \FunctionTok{theme\_minimal}\NormalTok{()}
\end{Highlighting}
\end{Shaded}

\pandocbounded{\includegraphics[keepaspectratio]{guide_files/figure-pdf/plot-compare-1.pdf}}

\section{Step 6: Evaluate Your
Results}\label{step-6-evaluate-your-results}

Once you have a curve and AUC, ask yourself:

\begin{itemize}
\tightlist
\item
  Does the curve bow toward the \textbf{upper-left} corner? If it hugs
  the diagonal, the test has little discriminative ability.
\item
  Is the \textbf{AUC above 0.5}? If not, check whether the
  \texttt{direction} argument is set correctly.
\item
  Does the \textbf{observation count} in the \texttt{pROC} output match
  your dataset? If not, check for unexpected missing values.
\item
  Is the \textbf{confidence interval} for the AUC reasonably narrow?
  Wide intervals suggest insufficient sample size.
\item
  Does the AUC make sense given what is known about this test from prior
  literature?
\end{itemize}

\section{Additional Resources}\label{additional-resources}

\begin{itemize}
\tightlist
\item
  \textbf{\texttt{pROC} vignette}: Run \texttt{vignette("pROC")} in R or
  see Robin et al. (2011) for the full paper
\item
  \textbf{Fawcett (2006)}: A readable conceptual introduction to ROC
  analysis (Fawcett 2006)
\item
  \textbf{Pepe (2003)}: A comprehensive textbook treatment of diagnostic
  test evaluation (\textbf{pepe2003?})
\end{itemize}

\bookmarksetup{startatroot}

\chapter*{References}\label{references}
\addcontentsline{toc}{chapter}{References}

\markboth{References}{References}

\phantomsection\label{refs}
\begin{CSLReferences}{1}{0}
\bibitem[\citeproctext]{ref-fawcett2006}
Fawcett, Tom. 2006. {``An Introduction to {ROC} Analysis.''}
\emph{Pattern Recognition Letters} 27 (8): 861--74.
\url{https://doi.org/10.1016/j.patrec.2005.10.010}.

\bibitem[\citeproctext]{ref-robin2011}
Robin, Xavier, Natacha Turck, Alexandre Hainard, Natalia Tiberti,
Frédérique Lisacek, Jean-Charles Sanchez, and Markus Müller. 2011.
{``{pROC}: An Open-Source Package for {R} and {S}+ to Analyze and
Compare {ROC} Curves.''} \emph{BMC Bioinformatics} 12 (1): 77.
\url{https://doi.org/10.1186/1471-2105-12-77}.

\end{CSLReferences}

\cleardoublepage
\phantomsection
\addcontentsline{toc}{part}{Appendices}
\appendix

\chapter{Topic}\label{topic}

Pick some subset of the information covered in your annual review
article that you are interested in exploring in more detail.

\begin{quote}
This user guide will focus on evaluating continuous diagnostic tests
using receiver operating characteristic or ROC curves. The guide will
cover how to generate the ROC curves from the given continuous test
results, how to interpret their shape, and how to summarize diagnostic
accuracy using the area under the curve, also referred to as AUC in the
article. This guide emphasizes practical implementation in R and assumes
users are familiar with basic statistical concepts such as sensitivity
and specificity, both essential the the interpretation of the ROC
curves.
\end{quote}

\begin{quote}
The guide will demonstrate how to perform ROC analysis using commonly
used R packages, such as pROC, including how to compute ROC curves and
AUC values and how to visualize results. Guidance will be provided on
when ROC curves are an appropriate tool for evaluating diagnostic
performance and when they may be misleading or insufficient, such as in
settings where disease prevalence or decision costs are critical, as
well as a guide for how to determine when using the ROC is an
appropriate analysis. This guide will also provide guidance on when this
analysis is acceptable, eg. for continous tests and not for binary
tests, and what format the data must be in in order to preform said
analysis.
\end{quote}

\begin{quote}
This user guide will not cover advanced topics discussed in the review
article, including covariate-adjusted ROC curves, time-dependent ROC
analysis, or optimal threshold selection methods such as the Youden
Index. These topics are excluded to maintain a clear focus on
foundational ROC analysis for continuous diagnostic tests.
\end{quote}

Identify the scope of your user guide - what will you cover? Specific
software packages? When to use or not to use a specific technique?

\begin{tcolorbox}[enhanced jigsaw, opacityback=0, opacitybacktitle=0.6, toptitle=1mm, leftrule=.75mm, colback=white, titlerule=0mm, title=\textcolor{quarto-callout-note-color}{\faInfo}\hspace{0.5em}{Dr.~V's comments}, coltitle=black, colbacktitle=quarto-callout-note-color!10!white, left=2mm, bottomrule=.15mm, bottomtitle=1mm, arc=.35mm, toprule=.15mm, rightrule=.15mm, breakable, colframe=quarto-callout-note-color-frame]

I think it would be helpful to set yourself up with a couple of
different real-world datasets (that are different from the data sets
used in the vignettes) for the \texttt{pROC} package. Ideally, you can
find one that shows off when ROC curves don't work well, and one that
shows when they do.

\end{tcolorbox}

\chapter{Needs Assessment}\label{needs-assessment}

What does someone trying to accomplish your chosen task need help with?

They would need the most help with the implementation of coding.
Assuming this guide is intended for applied researchers in biology and
medicine. Users new to using these packages or coding will likely
struggle a bit to connect the underlying concepts to actual
implementation in R.

What parts are likely to be tricky?

The trickiest part for most users is choosing the right R package and
functions. Several options exist (\texttt{pROC}, \texttt{ROCR},
\texttt{OptimalCutpoints}), and it is not always obvious which to use or
how the key functions relate to one another. Users may also be unsure
when ROC analysis is actually appropriate versus when it might be
misleading. Might also add some requirements for the data that is being
used.

What resources are already available on this topic that may be helpful?
Look for e.g.~software vignettes, package documentation, papers about
software packages, and so on.

Some that I found were: \texttt{pROC} package vignette Robin et
al.~(2011) paper in \emph{BMC Bioinformatics} Fawcett (2006)
\emph{Pattern Recognition Letters}

(will use these to help describe in more depth what the curves purposes
are and how to draw conclusions that are statistically correct based on
the produced results.)

\chapter{Needs Assessment}\label{needs-assessment-1}

\chapter{Task Analysis}\label{task-analysis}

Here are some questions to guide you through the process of doing a task
analysis. You don't have to specifically answer each one of these (and
some may not apply), but they should get you started thinking in the
right direction.

What are the prerequisites, for both knowledge and e.g.~software setup?

Users should understand basic diagnostic testing concepts: sensitivity,
specificity, and binary disease outcomes. Basic R skills are required
loading data, installing packages, and running functions. Users will
need R, RStudio, and the \texttt{pROC} and \texttt{ggplot2} packages
installed.

What are the basic requirements for any data that the method is used on?

\begin{itemize}
\item
  Are missing values allowed?
\item
  Should the data be confined within a certain range?
\item
  Does the data have to be approximately normally distributed?

  The dataset must have a numeric column for the continuous test result
  and a binary column for disease status (0/1 or a two-level factor).
  (ROC analysis is not appropriate for binary test results) Missing
  values are dropped by default in \texttt{pROC}, so users should check
  how many observations are affected. No distribution assumptions are
  required.
\end{itemize}

What are the basic components of the task? - Outline these in a bit more
detail

\begin{enumerate}
\def\labelenumi{\arabic{enumi}.}
\tightlist
\item
  Decide if ROC analysis is appropriate.(this includes identifieing if
  the results are continous or binary and if the use of a ROC curve is
  relevant)
\item
  Prepare the data. Confirm column types, check for missing values
\item
  Generate the ROC curve and AUC. Use \texttt{pROC::roc()} with the
  response and predictor variables, then \texttt{pROC::auc()} to get the
  AUC.
\item
  Visualize and interpret. Plot using \texttt{pROC::ggroc()} and
  annotate with the AUC.As describied in the article (basic
  interpretation), a curve bowing toward the upper-left indicates good
  discrimination; an AUC of 0.5 means no discriminative ability and 1.0
  means perfect.
\end{enumerate}

What decisions does the user need to be prepared to make on the fly?

\begin{itemize}
\tightlist
\item
  Is ROC analysis the right tool for this data and research question?
\item
  Is the test direction set correctly in \texttt{roc()}?
  (\texttt{direction\ =\ "\textless{}"} or
  \texttt{direction\ =\ "\textgreater{}"})
\item
  Is the AUC meaningful in this clinical or scientific context?
\end{itemize}

What questions should the user ask of the ``first draft'' of the
product? What adjustments may need to be made?

\begin{itemize}
\tightlist
\item
  Does the curve bow toward the upper-left as expected?
\item
  Does the observation count match your dataset?
\item
  Are results consistent with prior literature on this test?
\end{itemize}

adjustments

-If the AUC is below 0.5, the test direction is likely flipped and might
require user to reverse the direction argument.

\section{Additional Guidance}\label{additional-guidance}

Your check-in should answer these basic questions (and similar concerns
that apply more directly to your topic).

Once you've completed the check-in, you can use this section to
jump-start an introduction/set-up/getting started section in your user
guide. This document should remain as an appendix to your main report.




\end{document}
